\hypertarget{move-semantics-smart-pointers}{%
\chapter{MOVE SEMANTICS \& SMART
POINTERS}\label{move-semantics-smart-pointers}}

\hypertarget{rvalue-references-move-semantics}{%
\section{1. Rvalue References \& Move
Semantics}\label{rvalue-references-move-semantics}}

C++11 solves the problem of unnecessary copying with \textbf{Move
Semantics}.

\hypertarget{lvalues-vs-rvalues}{%
\subsection{1.1 Lvalues vs Rvalues}\label{lvalues-vs-rvalues}}

\begin{itemize}
\tightlist
\item
  \textbf{Lvalue (Left Value):} Has a name, has an address (identity).
  Persists beyond the expression.
\item
  \textbf{Rvalue (Right Value):} Temporary, no name (or about to die).
  Cannot take address.
\end{itemize}

\begin{lstlisting}[language={C++}]
int x = 10; // x is lvalue, 10 is rvalue
int y = x;  // y is lvalue
int z = x + y; // (x+y) is rvalue
\end{lstlisting}

\hypertarget{rvalue-references-t}{%
\subsection{\texorpdfstring{1.2 Rvalue References
(\texttt{T\&\&})}{1.2 Rvalue References (T\&\&)}}\label{rvalue-references-t}}

A reference that binds \emph{only} to rvalues. Represents an object we
can ``steal'' from.

\begin{lstlisting}[language={C++}]
void f(int& x)  { cout << "Lvalue ref"; }
void f(int&& x) { cout << "Rvalue ref"; }

int a = 5;
f(a); // Lvalue ref
f(5); // Rvalue ref
\end{lstlisting}

\hypertarget{move-constructor-move-assignment}{%
\subsection{1.3 Move Constructor \& Move
Assignment}\label{move-constructor-move-assignment}}

Instead of copying data (slow), we steal pointers (fast).

\begin{lstlisting}[language={C++}]
class Buffer {
    int* data;
    size_t size;
public:
    // Move Constructor
    Buffer(Buffer&& other) noexcept : data(other.data), size(other.size) {
        other.data = nullptr; // Nullify source to prevent double free
        other.size = 0;
    }

    // Move Assignment
    Buffer& operator=(Buffer&& other) noexcept {
        if (this != &other) {
            delete[] data;       // Free own resources
            data = other.data;   // Steal resources
            size = other.size;
            other.data = nullptr;// Nullify source
            other.size = 0;
        }
        return *this;
    }
};
\end{lstlisting}

\hypertarget{stdmove}{%
\subsection{\texorpdfstring{1.4
\texttt{std::move}}{1.4 std::move}}\label{stdmove}}

Casts an lvalue to an rvalue, enabling move semantics.

\begin{lstlisting}[language={C++}]
Buffer b1;
Buffer b2 = std::move(b1); // Calls Move Constructor
// b1 is now in valid but unspecified state (empty)
\end{lstlisting}

\textbf{Warning:} Do not use \passthrough{\lstinline!b1!} after moving
from it.

\begin{center}\rule{0.5\linewidth}{0.5pt}\end{center}

\hypertarget{smart-pointers-raii}{%
\section{2. Smart Pointers (RAII)}\label{smart-pointers-raii}}

Manual \passthrough{\lstinline!new!}/\passthrough{\lstinline!delete!} is
prone to leaks. C++11 introduces strict ownership models.

\hypertarget{stdunique_ptr}{%
\subsection{\texorpdfstring{2.1
\texttt{std::unique\_ptr}}{2.1 std::unique\_ptr}}\label{stdunique_ptr}}

\begin{itemize}
\tightlist
\item
  \textbf{Ownership:} Exclusive. Only one pointer owns the object.
\item
  \textbf{Copying:} Disabled (\passthrough{\lstinline!delete!}).
\item
  \textbf{Moving:} Allowed.
\item
  \textbf{Overhead:} Zero (same as raw pointer).
\end{itemize}

\begin{lstlisting}[language={C++}]
#include <memory>

std::unique_ptr<int> p1(new int(5));
// std::unique_ptr<int> p2 = p1; // ERROR: Cannot copy
std::unique_ptr<int> p3 = std::move(p1); // OK: Ownership transferred
\end{lstlisting}

\hypertarget{stdshared_ptr}{%
\subsection{\texorpdfstring{2.2
\texttt{std::shared\_ptr}}{2.2 std::shared\_ptr}}\label{stdshared_ptr}}

\begin{itemize}
\tightlist
\item
  \textbf{Ownership:} Shared (Reference Counted).
\item
  \textbf{Destruction:} Object deleted when last
  \passthrough{\lstinline!shared\_ptr!} dies.
\item
  \textbf{Overhead:} Control block on heap (ref counts).
\end{itemize}

\begin{lstlisting}[language={C++}]
auto sp1 = std::make_shared<int>(10); // Efficient allocation
auto sp2 = sp1; // Ref count = 2
\end{lstlisting}

\hypertarget{stdweak_ptr}{%
\subsection{\texorpdfstring{2.3
\texttt{std::weak\_ptr}}{2.3 std::weak\_ptr}}\label{stdweak_ptr}}

\begin{itemize}
\tightlist
\item
  \textbf{Ownership:} Non-owning observer of
  \passthrough{\lstinline!shared\_ptr!}.
\item
  \textbf{Use Case:} Break cyclic references (A-\textgreater B,
  B-\textgreater A).
\end{itemize}

\begin{lstlisting}[language={C++}]
std::shared_ptr<int> sp = std::make_shared<int>(42);
std::weak_ptr<int> wp = sp;

if (auto locked = wp.lock()) { // Check if object exists
    std::cout << *locked;
}
\end{lstlisting}

\hypertarget{stdmake_shared-vs-new}{%
\subsection{\texorpdfstring{2.4 \texttt{std::make\_shared} vs
\texttt{new}}{2.4 std::make\_shared vs new}}\label{stdmake_shared-vs-new}}

Always use \passthrough{\lstinline!std::make\_shared!}. It allocates the
object and control block in a single heap allocation (cache efficient).

\begin{center}\rule{0.5\linewidth}{0.5pt}\end{center}

\hypertarget{perfect-forwarding}{%
\section{3. Perfect Forwarding}\label{perfect-forwarding}}

Used in templates to preserve value category (lvalue vs rvalue).

\begin{lstlisting}[language={C++}]
template<typename T>
void wrapper(T&& arg) {
    // std::forward passes arg as lvalue if given lvalue,
    // rvalue if given rvalue.
    func(std::forward<T>(arg)); 
}
\end{lstlisting}
